
\paragraph{Further equilibrium results.} Four results in Supplement~\ref{app:eqinv} extend this section. None is needed for the main argument.
\begin{itemize}[itemsep=1pt,topsep=2pt]
\item \emph{A strategic insider} with short-lived information (Theorem~\ref{thm:insider}). The equilibrium is characterized in closed form and nests Kyle's one-period values. Its coefficients depend on the cohorts only through $G$, and the observed price is the closed loop driven through a known invertible moving-average filter. So the identification boundary and the $(\Delta\eta)^{-(4K-2)}$ rate hold with an insider present. Positive feedback raises total impact, and the insider then trades less and prices reveal less of her news.
\item \emph{Fragility} (Theorem~\ref{thm:fragility}). With positive feedback, a positive-impact equilibrium exists only if dealer risk aversion lies below a threshold. For strong feedback the threshold is inversely proportional to the feedback intensity $G(1)$. With several cohorts the bound depends on the population only through $G(1)$, again an aggregate.
\item \emph{Rents} (Theorem~\ref{thm:rents}). A dealer that misjudges the aggregate response by $\varepsilon$ hands an insider with a committed schedule a rent of exact order $\varepsilon$, or charges it one. Public data reveal the aggregate response at the ordinary rate. So the information that protects price formation is cheap to learn.
\item \emph{A targeted halt in equilibrium} (Theorem~\ref{thm:eqhalt}). Two markets that agree in equilibrium to within the cube of the memory spacing can differ at order one after their largest group is halted.
\end{itemize}
Disclosure keeps its value in equilibrium too (Corollary~\ref{cor:eqdisclosure}).
